\section{Conclusion}
\label{sec:conclusion}
%This paper introduces transition certificates for verifying LTL properties of discrete-time dynamical systems. Transition certificates reduce conservatism compared to state-triplet methods~\cite{wongpiromsarn2015automata} and have search spaces no larger than $\mathcal{X} \times Q \times Q$, significantly smaller than closure certificates~\cite{murali2024closure}. We present synthesis methods and demonstrate a speedup over closure certificates in several cases. Future work will focus on investigating its application to controller synthesis and adapting these notions for the verification and synthesis for continuous time systems.

In this paper, we have proposed transition certificates as an efficient framework for verifying Linear Temporal Logic (LTL) properties of discrete-time dynamical systems with continuous state spaces. Our approach leverages complementary certificates—transition safety and persistence certificates—to reduce conservatism compared to state-triplet methods and minimize search spaces relative to closure certificates, achieving significant speedups in LTL verification. We have developed automated synthesis methods based on polynomial and neural network templates, demonstrating practicality through case studies like the Kuramoto oscillator and two-room temperature model. Nonetheless, our method faces a few shortcomings that motivate future work. First, the synthesis process, particularly for neural network templates, can be computationally demanding for high-dimensional systems with highly nonlinear dynamics, due to the bottleneck in SMT solving. Second, the framework relies on the construction of NBA for LTL specifications, which involves exponential complexity and may limit scalability to complex formulas. Additionally, the current work focuses on discrete-time systems, leaving continuous-time systems as an open challenge. Future work will explore several directions. We plan to investigate the application of transition certificates to controller synthesis, enabling safe control policies with formal guarantees. Another direction is to adapt the notions for continuous-time systems, extending the verification framework to broader classes of dynamics. We also aim to integrate learning techniques, such as inferring specifications from trajectory data or automating template selection, to enhance adaptability and reduce manual effort. The transition certificates framework has potential applications beyond LTL verification, including the analysis of cyber-physical systems, AI-based controllers, and other domains requiring temporal logic guarantees. For instance, it could be used to ensure safety in autonomous systems or verify protocols in distributed computing. By providing a less conservative and efficient verification method, our work contributes to reliable system design, with implications for safety-critical areas like robotics and smart infrastructure. We will further explore these directions in future research.